\subsection{Incorporating ARMA-Noise into EM Algorithm}

With the analysis of ARMA-Noise processes already established \citep{WongMiller1990}, we examine how such an analysis may be applied, not to explicit realizations of ARMA-Noise processes, but rather estimates of one in the form of residuals from a fixed effects model. We define our model similarly to the ARE model, but with the latent time series $\uu$ now specified as an ARMA(1,1) model as opposed to strictly AR(1):

\begin{equation}
    \label{eq:fullmodel1}
    y_{ti} = \mathbf{x}_{ti}\beta + u_t + \epsilon_{ti}, \hspace{.4in} i = 1,...,n \hspace{.4 in} \epsilon_{ti} \sim \text{i.i.d. } N(0,\sigma_\epsilon^2) 
\end{equation}

\begin{equation}
    \label{eq:fullmodel2}
    u_t - \phi u_{t-1} = \eta_t + \theta \eta_{t-1} , \hspace{.4in} t=1,...T \hspace{.4in}  \hspace{.2in} \eta_t \sim N(0,\sigma_\eta^2) 
\end{equation}

The algorithm we propose is an iterative process applying the estimates of parameters from one level to update those of the other. The process begins similarly to that outlined by Modugno \textit{et al}, with a few key differences:

\begin{itemize}
    \item[1)] Estimate ${\beta}$ with $\hat{\beta} = (X_f'X_f)^{-1}X_f'\yy$;
    \item[2)] Estimate $X_t\uu + \varepsilon$ with $\yy - X_f\hat{\beta}$ and fit ARMA-Noise model to obtain estimates $(\hat{\phi}, \hat{\theta}, \hat{\sigma_\eta^2}, \hat{\sigma_\varepsilon^2})$;
    \item[3)] Calculate $\hat{\Omega}$ and $l(\hat{\xi})$;
    \item[4)] Refit fixed effect parameter estimates by calculating GLS estimate $\hat{\beta}_{GLS} = (X_f'\hat{\Omega}^{-1}X_f)X_f'\hat{\Omega}^{-1}\yy$;
    \item[5)] Repeat Steps 2-4 until convergence criteria are met.;
    \item[6)] Generate confidence regions and eliminate insignificant parameters if applicable.
\end{itemize}

\subsection{Evaluating Convergence}


Instead of splitting the likelihood by conditioning on $\hat{\uu}$, doing so with $\bar{\yy}$ instead of has the advantage that likelihoods are based on a constant statistic, as opposed to the ever-changing estimate of $\uu$.

Let's recall the log likelihood $l(\xi)$:

$$ l(\xi) = -\frac{N}{2}\text{ln}(2\pi) - \frac{1}{2}\text{ln}|\Omega| - \frac{1}{2}(\mathbf{y}-X_f\beta)'\Omega^{-1}(\mathbf{y} - X_f\beta) $$

We now condition on $\bar{\mathbf{y}}_{t\cdot} = n^{-1}X_t'\mathbf{y}$, which has the following distribution

$$\bar{\mathbf{y}} \sim N(n^{-1}X_t'X_f\beta, n^{-2}X_t'\Omega X_t) $$
After simplifying, we get the following expression for $\Sigma_{\bar{\mathbf{y}}}$:

$$ \Sigma_{\bar{\mathbf{y}}} = \Gamma + n^{-1}\sigma_\varepsilon^2I_T, $$
which has the following log likelihood:

$$ l_t(\xi) = \text{ln}[f(\bar{\mathbf{y}}|\xi)] = -\frac{T}{2}\text{ln}(2\pi) - \frac{1}{2}\text{ln}|\Sigma_{\bar{\mathbf{y}}}| - \frac{1}{2}(\bar{\mathbf{y}} - n^{-1}X_t'X_f\beta)'\Sigma_{\bar{\mathbf{y}}}^{-1}(\bar{\mathbf{y}} - n^{-1}X_t'X_f\beta) $$
If we then condition $\mathbf{y}$ on $\bar{\mathbf{y}}$, the only randomness left is in the vector of deviations, which are independent of $\bar{\mathbf{y}}$:

$$ \mathbf{y} - X_t\bar{\mathbf{y}} = \bm{\varepsilon} - X_t\mathbf{a}   $$

Let's establish the distribution of this vector of deviations. The mean is immediately zero, and the covariance is given by:

$$ \text{Var}(\bm{\varepsilon} - X_t\mathbf{a}) = \text{Var}(\bm{\varepsilon})  - 2 \text{Cov}(\bm{\varepsilon},X_t\mathbf{a}) + \text{Var}(X_t\mathbf{a})  $$
% Remebering $\mathbf{a} = n^{-1}X_t'\bm{\varepsilon}$ the above expression becomes:

% $$ \text{Cov}(\bm{\varepsilon} - X_t\mathbf{a}) = \text{Var}(\bm{\varepsilon})  - 2 \text{Cov}(\bm{\varepsilon},n^{-1}X_tX_t'\bm{\varepsilon}) + \text{Var}(n^{-1}X_tX_t'\bm{\varepsilon})  $$,
which reduces to :

$$ \text{Var}(\bm{\varepsilon} - X_t\mathbf{a}) = \sigma_\varepsilon^2I_N  - 2 \sigma_\varepsilon^2 n^{-1}X_tX_t' + n^{-2}\sigma_\varepsilon^2(X_tX_t')^2  $$
Because $X_t'X_t = nI_N$, we get:
\begin{equation}
    \text{Var}(\bm{\varepsilon} - X_t\mathbf{a}) = \sigma_\varepsilon^2(I_N  -  n^{-1}X_tX_t') 
    \label{eq:var_dev}
\end{equation}


We now derive the second part of our split likelihood. Because of the rank-deficient nature of the matrix Equation \ref{eq:var_dev}, we can simplify the log likelihood

\begin{equation}
l_i(\xi) = \ln f(\mathbf{y} \mid \bar{\mathbf{y}}, \sigma_\varepsilon^2)
= -\frac{N-T}{2} \ln(2 \pi \sigma_\varepsilon^2)
- \frac{1}{2 \sigma_\varepsilon^2} \sum_{t=1}^{T} \sum_{i=1}^{n} 
\left( y_{it} - \bar{y}_t \right)^2,
\end{equation}


With this we can express the full likelihood as $l(\xi) = l_t(\xi) + l_i(\sigma_\varepsilon^2)$. However, since $\bar{\mathbf{y}}_{t\cdot}$ is independent of our model selection and parameter estimates, $l_i$ is going to be relatively stable. We therefore pair $\xi \in \mathbb{R}^k$ and $l_t(\xi)$ to calculate the Akaike Information Criterion (AIC): 

$$ AIC = -2k -2l_t(\xi) $$

Calculating this weighted likelihood is how we will compare the performance of models that include different subsets of parameters.




\begin{figure}[htbp]
    \centering
    \vspace{.5in}
\resizebox{\textwidth}{!}{
\begin{tikzpicture}[scale=1.2, every node/.style={transform shape}, node distance=2.5cm]{}
\small
    % Nodes vertically aligned
    \node (start) [startstop] {$\mathbf{y} = X_f\beta + X_r\mathbf{b} + X_t\mathbf{u} + \varepsilon$};
    \node (step1) [process, right=1.5cm of start] 
        {Fit ME Model to $\yy$};
    \node (step2) [process, below = 1cm of start] 
        {$\mathbf{z}:= \yy - X_f\hat{\beta} - X_r\hat{\mathbf{b}}$};
    \node (step3) [process, below = 2cm of step2] 
        {${\mathbf{w}}_t = n^{-1}\sum_{i=1}^{n} \mathbf{z}_{it}$};
    
    % Parallel branch
    \node (step4a) [process, below = 4cm of step1] 
        {${\mathbf{x}} =  \mathbf{z} - X_t\mathbf{w}$};
    \node (step5a) [process, below = 2cm of step4a] 
        {$\hat{\sigma}_a^{2} = \frac{||\mathbf{x}||^2\sum n_t^{-1}}{(N-T)T}$};
    
    \node (step4b) [process, below = 2cm of step3] 
        {$(\hat{\theta}_*, \hat{\phi}, \hat{\sigma}_*^2)$};
    
    \node (step5b) [process, below  = 2cm of step5a] 
        {$H(\theta,\sigma_\eta^2)$};
    
    \node (step7) [process, above right = 1cm and 1cm of step5b] 
        {$(\hat{\theta}, \hat{\sigma}_\eta^2)$};
    
    % % Merge after parallel
    \node (step8) [process, above = 2cm of step7] 
        {Calculate $\hat{\beta}_{GLS}^{(k)}$ and $\mathbb{E}(\mathbf{b}|\mathbf{y})$};
    
    \node (checkll) [decision, above = 2cm of step8] 
        {Likelihood converged?};
    \node (done) [startstop, above = 2cm of checkll] {Done};
    
    % Arrows
    \draw [arrow] (start) -- (step1)
        node[pos=0.5, anchor = south]{$k = 1$};
    \draw [arrow] (step1) |- (step2) node[pos=0.5, anchor = south]{};
    \draw [arrow] (step2) -- (step3)
        node[pos=.5,anchor = east]{Calculate}
        node[pos=.5,anchor=west]{Daily Averages};
    \draw [arrow] (step3) -- (step4a) 
        node[pos=.5,anchor = south]{}
        node[pos=.5,anchor=north]{};
    \draw [arrow] (step3) -- (step4b) 
        node[pos=0.5, anchor = east]{(fit ARMA}
        node[pos=.5,anchor=west]{to $\mathbf{w}$)};
    \draw [arrow] (step4a) -- (step5a);
    \draw [arrow] (step5a) -- (step5b);
    \draw [arrow] (step4b) |- (step5b);
    \draw [arrow] (step5b) -| (step7)
        node[pos=.6,anchor = east]{Solve for}
        node[pos=.6,anchor = west]{roots};
    \draw [arrow] (step7) -- (step8)
        node[pos=.5, anchor = east]{Update}
        node[pos=.5, anchor = west]{$(\hat{\GG},\hat{\Omega})$};
    \draw [arrow] (step8) -- (checkll)
        node[pos=0.5, right]{$k = k+1$};
    \draw [arrow] (checkll) -- (done)
        node[pos=.75, anchor = west]{Yes};
    \draw [arrow] (checkll) |- (step2) 
        node[pos=.6, anchor = south]{No};
\end{tikzpicture}
}
    \caption{Flowchart of algorithm steps for estimating model parameters}
    \label{fig:modelflow}
\end{figure}